\sectionheading{5.49 বহুনির্বাচনি প্রশ্ন (MCQ)}
\begin{infobox}
\textbf{কাঠামো:} প্রথমে Topic-wise Practice—নম্বর ১ থেকে ধারাবাহিক। এরপর Board-style Mixed MCQ—আবার ১ থেকে শুরু। গুরুত্বপূর্ণ Topic-এ একই Concept-এর ভিন্ন angle থেকে একাধিক প্রশ্ন রাখা হয়েছে; অপ্রয়োজনীয় repetition রাখা হয়নি। শেষে নির্বাচিত Difficult, Critical ও Trap প্রশ্নের ব্যাখ্যা এবং একই Concept থেকে আর কীভাবে প্রশ্ন হতে পারে তা দেওয়া হয়েছে।
\end{infobox}

\partheading{অংশ–১ : Topic-wise MCQ Practice}
\begin{multicols}{2}
\setlength{\columnsep}{4mm}
\raggedcolumns
\topicheading{A. Program, Programming Language ও ভাষার স্তর}
\mcq
{1}
{কম্পিউটারকে নির্দিষ্ট কাজ করানোর জন্য সাজানো নির্দেশনার সমষ্টিকে কী বলে?}
{Algorithm}
{Flowchart}
{Program}
{Translator}

\mcq
{2}
{Program লেখার জন্য ব্যবহৃত নিয়ম ও প্রতীকের সমষ্টিকে কী বলা হয়?}
{Hardware}
{Operating System}
{Database}
{Programming Language}

\mcq
{3}
{কম্পিউটারের সরাসরি বোধগম্য ভাষা কোনটি?}
{Python}
{High Level Language}
{Machine Language}
{Assembly Language}

\mcq
{4}
{Machine Language প্রধানত কোন দুটি Digit দিয়ে গঠিত?}
{0 ও 9}
{1 ও 2}
{A ও B}
{0 ও 1}

\mcq
{5}
{প্রথম প্রজন্মের Programming Language কোনটি?}
{High Level Language}
{Machine Language}
{Assembly Language}
{4GL}

\mcq
{6}
{দ্বিতীয় প্রজন্মের Programming Language কোনটি?}
{Machine Language}
{4GL}
{Python}
{Assembly Language}

\mcq
{7}
{Assembly Language-এ Instruction লেখার জন্য সাধারণত কী ব্যবহৃত হয়?}
{Flowchart Symbol}
{শুধু Binary Digit}
{Natural Sentence}
{Mnemonic Code}

\mcq
{8}
{Machine Language ও Assembly Language-কে একত্রে সাধারণত কী বলা হয়?}
{High Level Language}
{Low Level Language}
{Visual Language}
{4GL}

\mcq
{9}
{নিচের কোনটি Middle Level Language-এর প্রচলিত উদাহরণ?}
{Machine Code}
{C}
{Assembly Code}
{SQL}

\mcqwide
{10}
{High Level Language-এর প্রধান সুবিধা কোনটি?}
{Assembler ছাড়া চলে না}
{Hardware-এর সাথে অত্যন্ত ঘনিষ্ঠ}
{মানুষের জন্য তুলনামূলক সহজ}
{শুধু 0 ও 1 দিয়ে লিখতে হয়}

\mcqwide
{11}
{Machine Language-এর একটি প্রধান বৈশিষ্ট্য কোনটি?}
{Machine-dependent}
{Mnemonic-based}
{সম্পূর্ণ machine-independent}
{Natural language-based}

\mcq
{12}
{Assembly Language কোনটির তুলনায় মানুষের জন্য তুলনামূলক সহজ?}
{Python}
{4GL}
{Machine Language}
{High Level Language}

\mcq
{13}
{নিচের কোনটি High Level Language?}
{Python}
{Assembly Language}
{Binary Code}
{Machine Language}

\mcq
{14}
{নিচের কোনটি High Level Language-এর উদাহরণ?}
{Machine Code}
{Mnemonic মাত্র}
{Java}
{Opcode}

\mcqwide
{15}
{Python সম্পর্কে সঠিক বক্তব্য কোনটি?}
{এটি High Level Language}
{এটি শুধু Hardware Language}
{এটি কেবল Assembler}
{এটি Machine Language}

\mcqwide
{16}
{4GL-এর পূর্ণরূপ কোনটি?}
{Fourth General Logic}
{Fourth Generation Language}
{Four General Language}
{Fast Generation Language}

\mcq
{17}
{4GL সাধারণত কোন উদ্দেশ্যে High Level Language-এর চেয়েও বেশি user-oriented?}
{কম Code-এ কাজ সহজ করা}
{শুধু CPU নিয়ন্ত্রণ}
{শুধু Binary লেখা}
{শুধু Mnemonic লেখা}

\mcqwide
{18}
{High Level Language-এর Code সাধারণত সরাসরি CPU বুঝতে পারে কি?}
{শুধু Integer থাকলে পারে}
{না, Translator প্রয়োজন হয়}
{শুধু Loop থাকলে পারে}
{হ্যাঁ, সবসময়}

\mcq
{19}
{Low Level Language-এর সাথে কোনটির সম্পর্ক বেশি ঘনিষ্ঠ?}
{Document Formatting}
{Natural Language}
{Web Browser}
{Hardware}

\mcqwide
{20}
{নিচের কোন বক্তব্যটি সঠিক?\newline i. Machine Language machine-dependent।\newline ii. Assembly Language-এ mnemonic ব্যবহৃত হয়।\newline iii. Python High Level Language।}
{i, ii ও iii}
{i ও iii}
{i ও ii}
{ii ও iii}

\topicheading{B. Translator Program, Compiler, Interpreter ও Assembler}
\mcq
{21}
{এক ভাষায় লেখা Program-কে অন্য ভাষায় রূপান্তরকারী Program-কে কী বলে?}
{Debugger}
{Translator Program}
{Editor}
{Loader}

\mcq
{22}
{পুরো Source Program একসঙ্গে অনুবাদ করে কোনটি?}
{Interpreter}
{Linker}
{Compiler}
{Assembler}

\mcq
{23}
{সাধারণত Statement ধরে ধরে অনুবাদ ও Execute করে কোনটি?}
{Interpreter}
{Editor}
{Compiler}
{Assembler}

\mcq
{24}
{Assembly Language-কে Machine Language-এ রূপান্তর করে কোনটি?}
{Compiler}
{Debugger}
{Interpreter}
{Assembler}

\mcq
{25}
{Programmer যে মূল Code লেখে তাকে কী বলা হয়?}
{Object Program}
{Output File}
{Machine Object}
{Source Program}

\mcq
{26}
{Translation-এর পরে প্রাপ্ত Machine-oriented Program-কে সাধারণত কী বলা হয়?}
{Algorithm}
{Flowchart}
{Source Program}
{Object Program}

\mcq
{27}
{Compiler-এর ক্ষেত্রে Error সাধারণত কখন দেখা যায়?}
{কখনোই নয়}
{Compilation-এর সময়}
{শুধু Input-এর সময়}
{Computer বন্ধ থাকলে}

\mcqwide
{28}
{Interpreter কোনো Statement-এ Error পেলে সাধারণত কী করে?}
{সবসময় Object Code বানায়}
{Assembler চালায়}
{সব Error উপেক্ষা করে}
{সেই Error ঠিক না হওয়া পর্যন্ত পরের Statement-এ যায় না}

\mcqwide
{29}
{একবার সফলভাবে Compile করা Program বারবার Execute করার ক্ষেত্রে কোনটির সুবিধা বেশি?}
{Compiler-ভিত্তিক Translation}
{কোনোটিই নয়}
{Assembler ছাড়া সব}
{প্রতিবার Interpreter দিয়ে line-by-line Translation}

\mcqwide
{30}
{Machine Language-এর Program চালাতে Translation প্রয়োজন হয় কি?}
{সাধারণত না}
{শুধু Interpreter লাগে}
{হ্যাঁ, Assembler লাগে}
{হ্যাঁ, অবশ্যই Compiler লাগে}

\mcq
{31}
{Assembler-এর Input কোন ভাষায় লেখা হয়?}
{High Level Language}
{Natural Language}
{Assembly Language}
{Machine Language}

\mcq
{32}
{Compiler ও Interpreter উভয়ই মূলত কোন শ্রেণির Program?}
{Device Driver}
{Operating System}
{Translator}
{Database}

\mcq
{33}
{High Level Language অনুবাদের জন্য কোনটি ব্যবহৃত হতে পারে?}
{Compiler বা Interpreter}
{শুধু Assembler}
{শুধু Editor}
{শুধু Loader}

\mcqwide
{34}
{নিচের কোন জোড়াটি সঠিক?}
{Interpreter—শুধু Flowchart আঁকে}
{Assembler—Assembly Language}
{Editor—Translator}
{Compiler—শুধু Machine Language লিখে}

\topicheading{C. Program Organization, Design, Development ও Computational Thinking}
\mcq
{35}
{Program-এর বিভিন্ন অংশকে পরিকল্পিতভাবে সাজানোর ধারণাকে কী বলা যায়?}
{Program Organization}
{Machine Code}
{File Compression}
{Data Entry}

\mcq
{36}
{Program Development-এর শুরুতে সবচেয়ে গুরুত্বপূর্ণ কাজ কোনটি?}
{File delete করা}
{Output print করা}
{Random Code লেখা}
{সমস্যা বোঝা ও বিশ্লেষণ}

\mcq
{37}
{Coding-এর আগে Logic পরিকল্পনার জন্য কোন দুটি উপায় বেশি ব্যবহৃত হয়?}
{Compiler ও Assembler}
{Input ও Output}
{File ও Folder}
{Algorithm ও Flowchart}

\mcq
{38}
{Program-এর ভুল শনাক্ত ও সংশোধনের প্রক্রিয়াকে কী বলা হয়?}
{Formatting}
{Installing}
{Debugging}
{Encrypting}

\mcq
{39}
{Program সঠিক ফল দিচ্ছে কি না যাচাই করাকে কী বলা হয়?}
{Scanning}
{Typing}
{Linking}
{Testing}

\mcq
{40}
{Program সম্পর্কিত প্রয়োজনীয় লিখিত তথ্য ও নির্দেশনা সংরক্ষণকে কী বলা হয়?}
{Iteration}
{Compilation}
{Documentation}
{Translation}

\mcq
{41}
{ব্যবহারের পরে Program সংশোধন ও উন্নয়ন করাকে কী বলা হয়?}
{Maintenance}
{Initialization}
{Input}
{Tokenization}

\mcq
{42}
{বড় Program-কে ছোট ছোট অংশে ভাগ করে তৈরির পদ্ধতি কোনটি?}
{Machine Programming}
{Direct Programming}
{Modular Programming}
{Binary Programming}

\mcqwide
{43}
{Modular Programming-এর সুবিধা কোনটি?}
{সব Variable Global হয়}
{Error কখনো হয় না}
{Algorithm প্রয়োজন হয় না}
{Program বোঝা ও রক্ষণাবেক্ষণ সহজ হয়}

\mcqwide
{44}
{Computational Thinking-এর মূল লক্ষ্য কী?}
{সমস্যাকে যৌক্তিকভাবে বিশ্লেষণ করে সমাধান তৈরি}
{শুধু Output লিখে দেওয়া}
{শুধু Code মুখস্থ করা}
{শুধু Hardware তৈরি}

\mcq
{45}
{সমস্যাকে ছোট ছোট অংশে ভাগ করে বোঝা কোন চিন্তাপদ্ধতির সাথে সম্পর্কিত?}
{Machine Translation}
{Computational Thinking}
{File Handling}
{Formatting}

\mcqwide
{46}
{Problem Analysis-এর সময় কোনটি আগে নির্ধারণ করা দরকার?}
{শুধু Font}
{শুধু Compiler Version}
{কী Input, কী Processing, কী Output}
{শুধু File Name}

\mcqwide
{47}
{একটি ভালো Program Design-এর বৈশিষ্ট্য কোনটি?}
{কোনো Test নেই}
{অপ্রয়োজনীয় repetition বেশি}
{সব কাজ এক লাইনে}
{Logic পরিষ্কার ও ধাপগুলো সংগঠিত}

\mcqwide
{48}
{Program Development-এর যুক্তিসংগত ক্রম কোনটি?}
{Output \(\rightarrow\) Coding \(\rightarrow\) Design \(\rightarrow\) Problem}
{Testing \(\rightarrow\) Coding \(\rightarrow\) Analysis \(\rightarrow\) Input}
{Analysis \(\rightarrow\) Design \(\rightarrow\) Coding \(\rightarrow\) Testing}
{Coding \(\rightarrow\) Analysis \(\rightarrow\) Output \(\rightarrow\) Problem}

\topicheading{D. IPO Model, Dry Run, Trace Table, Algorithm ও Flowchart}
\mcqwide
{49}
{IPO-এর পূর্ণরূপ কোনটি?}
{Input–Program–Operation}
{Input–Processing–Output}
{Input–Print–Output}
{Instruction–Process–Output}

\mcq
{50}
{IPO Model-এ ব্যবহারকারীর কাছ থেকে গ্রহণ করা Data কোন অংশ?}
{Output}
{Storage}
{Input}
{Processing}

\mcq
{51}
{IPO Model-এ Calculation বা Decision কোন অংশে পড়ে?}
{Processing}
{Input}
{Output}
{Comment}

\mcq
{52}
{Program-এর চূড়ান্ত ফল প্রদর্শন কোন অংশ?}
{Input}
{Output}
{Variable}
{Processing}

\mcq
{53}
{দৈর্ঘ্য ও প্রস্থ নিয়ে আয়তক্ষেত্রের ক্ষেত্রফল বের করার ক্ষেত্রে Processing কোনটি?}
{দৈর্ঘ্য × প্রস্থ}
{ক্ষেত্রফল প্রদর্শন}
{দৈর্ঘ্য ও প্রস্থ গ্রহণ}
{Program বন্ধ করা}

\mcq
{54}
{Program বাস্তবে Run না করে হাতে ধাপে ধাপে Logic পরীক্ষা করাকে কী বলে?}
{Installation}
{Compilation}
{Formatting}
{Dry Run}

\mcq
{55}
{Dry Run-এর সময় Variable-এর মান ধাপে ধাপে Table-এ লিখলে তাকে কী বলা হয়?}
{Symbol Table}
{Data Table}
{Trace Table}
{Truth Table}

\mcqwide
{56}
{Trace Table-এর প্রধান কাজ কোনটি?}
{Font পরিবর্তন করা}
{Code Translate করা}
{Variable-এর Value পরিবর্তন অনুসরণ করা}
{File Size কমানো}

\mcqwide
{57}
{Algorithm কী?}
{শুধু Diagram}
{সমস্যা সমাধানের সুনির্দিষ্ট ও সসীম ধাপসমূহ}
{Machine Code}
{শুধু Python Code}

\mcq
{58}
{একটি Algorithm-এর ধাপ কেমন হওয়া উচিত?}
{অসীম}
{শুধু Binary}
{স্পষ্ট ও দ্ব্যর্থহীন}
{অব্যবস্থিত}

\mcq
{59}
{Algorithm-এর গুরুত্বপূর্ণ বৈশিষ্ট্য কোনটি?}
{কখনো শেষ না হওয়া}
{শুধু English হওয়া}
{শুধু Diagram হওয়া}
{সসীম ধাপে শেষ হওয়া}

\mcqwide
{60}
{Algorithm সাধারণত Programming Language-এর উপর নির্ভরশীল কি?}
{না, মূল Logic ভাষা-স্বাধীন হতে পারে}
{হ্যাঁ, সম্পূর্ণ}
{শুধু C-এর উপর}
{শুধু Python-এর উপর}

\mcqwide
{61}
{Flowchart কী?}
{Algorithm-এর চিত্রভিত্তিক উপস্থাপন}
{Object Code}
{Database Table}
{Source Code}

\mcq
{62}
{Flowchart-এ Start/End বোঝাতে কোন Shape ব্যবহৃত হয়?}
{Parallelogram}
{Rectangle}
{Oval}
{Diamond}

\mcq
{63}
{Flowchart-এ Input/Output বোঝাতে কোন Shape ব্যবহৃত হয়?}
{Oval}
{Rectangle}
{Parallelogram}
{Diamond}

\mcq
{64}
{Flowchart-এ Processing বোঝাতে কোন Shape ব্যবহৃত হয়?}
{Diamond}
{Rectangle}
{Circle}
{Oval}

\mcq
{65}
{Flowchart-এ Decision বোঝাতে কোন Shape ব্যবহৃত হয়?}
{Rectangle}
{Oval}
{Parallelogram}
{Diamond}

\mcq
{66}
{Flowchart-এর এক ধাপ থেকে অন্য ধাপে প্রবাহ দেখাতে কী ব্যবহৃত হয়?}
{Variable}
{Arrow/Flow line}
{Keyword}
{Compiler}

\mcqwide
{67}
{Flowchart-এ Connector ব্যবহারের উদ্দেশ্য কী?}
{Variable Declare করা}
{Code Compile করা}
{দূরের Flow অংশকে সংযুক্ত করা}
{Output Print করা}

\mcqwide
{68}
{Algorithm ও Flowchart-এর সম্পর্ক কোনটি?}
{দুটিই Problem Logic প্রকাশ করে}
{দুটিই Translator}
{দুটিই Data Type}
{দুটিই Hardware}

\mcq
{69}
{Coding-এর আগে Algorithm/Flowchart তৈরির সুবিধা কী?}
{Input দরকার হয় না}
{Logic পরিষ্কার হয়}
{Error সবসময় অসম্ভব হয়}
{Compiler লাগে না}

\mcq
{70}
{নিচের কোনটি সঠিক মিল?}
{Decision—Diamond}
{Input—Rectangle}
{Processing—Oval}
{Start—Parallelogram}

\mcq
{71}
{একটি সংখ্যার জোড়/বিজোড় নির্ণয়ে Flowchart-এ কোন Symbol অবশ্যই দরকার?}
{Comment}
{Decision}
{Connector মাত্র}
{Terminal মাত্র}

\mcq
{72}
{Dry Run সবচেয়ে বেশি সাহায্য করে কোন Error ধরতে?}
{Power Error}
{Logical Error}
{Keyboard Error}
{Printer Error}

\topicheading{E. Python পরিচিতি, বৈশিষ্ট্য, Environment ও Syntax}
\mcq
{73}
{Python কোন ধরনের Programming Language?}
{শুধু 4GL}
{High Level}
{Machine}
{Assembly}

\mcqwide
{74}
{Python Program সাধারণত কীভাবে Execute হয়?}
{শুধু Machine Code হিসেবে}
{Interpreter-এর মাধ্যমে}
{Assembler-এর মাধ্যমে}
{Translator ছাড়াই সব ক্ষেত্রে}

\mcq
{75}
{Python কি Case-sensitive?}
{শুধু Number-এর ক্ষেত্রে}
{শুধু Comment-এর ক্ষেত্রে}
{হ্যাঁ}
{না}

\mcq
{76}
{Python-এ Code Block নির্ধারণে কোনটি বিশেষ গুরুত্বপূর্ণ?}
{Binary Prefix}
{Indentation}
{Semicolon বাধ্যতামূলক}
{Line Number}

\mcq
{77}
{Python Source File-এর প্রচলিত Extension কোনটি?}
{.exe}
{.java}
{.asm}
{.py}

\mcq
{78}
{Python-এর Interactive Mode-এ সাধারণত কোন Prompt দেখা যায়?}
{\%\%\%}
{\#\#\#}
{>>>}
{\$\$\$}

\mcq
{79}
{একাধিক Line-এর Program File-এ লিখে চালানোকে কোন Mode-এর সাথে বেশি সম্পর্কিত?}
{Script Mode}
{Assembly Mode}
{Binary Mode}
{Machine Mode}

\mcqwide
{80}
{IDLE কী?}
{Python-এর একটি Development Environment}
{একটি Operator}
{একটি Loop}
{একটি Data Type}

\mcq
{81}
{Python-এ Single-line Comment সাধারণত কোন Symbol দিয়ে শুরু হয়?}
{\#}
{//}
{--}
{/*}

\mcq
{82}
{Python Keyword কি Variable Name হিসেবে ব্যবহার করা যায়?}
{হ্যাঁ, সবসময়}
{শুধু Loop-এর ভিতরে}
{না}
{শুধু int হলে}

\mcq
{83}
{নিচের কোনটি বৈধ Python Identifier?}
{2marks}
{class}
{total\_marks}
{total-marks}

\mcq
{84}
{Python Identifier-এর প্রথম Character কী হতে পারে?}
{Letter বা underscore}
{শুধু space}
{শুধু digit}
{শুধু hyphen}

\mcq
{85}
{Python Identifier-এর মধ্যে Space ব্যবহার করা যায় কি?}
{হ্যাঁ}
{শুধু শেষে}
{না}
{শুধু শুরুতে}

\mcq
{86}
{if, for, while, def-এর Header-এর শেষে সাধারণত কোন Symbol থাকে?}
{;}
{,}
{:}
{.}

\mcq
{87}
{Python-এর Syntax Error-এর একটি সাধারণ কারণ কোনটি?}
{সঠিক Variable Name}
{সঠিক Colon}
{সঠিক Operator}
{ভুল Indentation}

\mcqwide
{88}
{Python-এর জনপ্রিয়তার একটি কারণ কোনটি?}
{শুধু Binary Code}
{সহজ ও পাঠযোগ্য Syntax}
{কেবল একটি Data Type}
{Assembler বাধ্যতামূলক}

\topicheading{F. Variable, Data Type, Input/Output ও Operator}
\mcq
{89}
{Python-এ Variable-এ Value দেওয়ার জন্য কোন Operator ব্যবহৃত হয়?}
{>=}
{!=}
{==}
{=}

\mcq
{90}
{Python-এ একই Variable পরে অন্য ধরনের Value ধারণ করতে পারে—এটি কোন বৈশিষ্ট্যের সাথে সম্পর্কিত?}
{Static Flowchart}
{Assembler}
{Machine Code}
{Dynamic Typing}

\mcq
{91}
{নিচের কোনটি বৈধ Variable Name?}
{score1}
{for}
{1score}
{score-1}

\mcq
{92}
{Python-এ \py{age} এবং \py{Age} কি একই Variable?}
{শুধু Loop-এ}
{হ্যাঁ}
{না}
{শুধু String হলে}

\mcq
{93}
{পূর্ণসংখ্যার Data Type কোনটি?}
{int}
{bool}
{float}
{str}

\mcq
{94}
{দশমিক সংখ্যার Data Type কোনটি?}
{float}
{int}
{str}
{tuple}

\mcq
{95}
{জটিল সংখ্যার Data Type কোনটি?}
{set}
{complex}
{bool}
{dict}

\mcq
{96}
{True ও False কোন Data Type-এর Value?}
{int}
{list}
{str}
{bool}

\mcq
{97}
{Text Data সংরক্ষণের জন্য কোন Data Type ব্যবহৃত হয়?}
{float}
{int}
{bool}
{str}

\mcq
{98}
{Ordered এবং Mutable Collection কোনটি?}
{set}
{tuple}
{NoneType}
{list}

\mcq
{99}
{Ordered কিন্তু Immutable Collection কোনটি?}
{dict}
{set}
{tuple}
{list}

\mcq
{100}
{Unique Element সংরক্ষণকারী Collection কোনটি?}
{tuple}
{set}
{str}
{list}

\mcq
{101}
{Key–Value Pair আকারে Data সংরক্ষণ করে কোনটি?}
{tuple}
{set}
{dict}
{list}

\mcq
{102}
{কোনো Value অনুপস্থিত বা নেই বোঝাতে Python-এ কোনটি ব্যবহৃত হয়?}
{empty}
{None}
{0 মাত্র}
{False মাত্র}

\mcq
{103}
{কোনো Value-এর Data Type জানার জন্য কোন Function ব্যবহার করা হয়?}
{type()}
{range()}
{print()}
{input()}

\mcq
{104}
{\py{input()} সাধারণত কোন Data Type Return করে?}
{str}
{int}
{bool}
{float}

\mcq
{105}
{User-এর Input-কে Integer-এ রূপান্তরের জন্য কোনটি ব্যবহার করা হয়?}
{bool()}
{int()}
{str()}
{type()}

\mcq
{106}
{User-এর Input-কে Decimal Number-এ রূপান্তরের জন্য কোনটি ব্যবহার করা হয়?}
{list()}
{float()}
{set()}
{int()}

\mcq
{107}
{Screen-এ Output দেখাতে কোন Function ব্যবহৃত হয়?}
{input()}
{range()}
{type()}
{print()}

\mcq
{108}
{\py{print()}-এ একাধিক Item-এর মাঝখানের Separator পরিবর্তনে কোন Argument ব্যবহৃত হয়?}
{end}
{step}
{sep}
{start}

\mcq
{109}
{\py{print()}-এর শেষে কী থাকবে তা নির্ধারণে কোন Argument ব্যবহৃত হয়?}
{stop}
{next}
{end}
{sep}

\mcq
{110}
{Variable-এর Value সহ সুন্দর String Output তৈরির আধুনিক উপায় কোনটি?}
{Assembler String}
{Flow String}
{Binary String}
{f-string}

\mcq
{111}
{নিচের কোনটি Arithmetic Operator?}
{==}
{and}
{in}
{+}

\mcqwide
{112}
{\textenglish{\ttfamily 10 / 4}-এর ফল কত?}
{0}
{2.5}
{4}
{2}

\mcqwide
{113}
{\textenglish{\ttfamily 10 // 4}-এর ফল কত?}
{2}
{2.5}
{0}
{4}

\mcqwide
{114}
{\textenglish{\ttfamily 10 \char37 4}-এর ফল কত?}
{2.5}
{2}
{4}
{1}

\mcqwide
{115}
{\textenglish{\ttfamily 2 ** 3}-এর ফল কত?}
{9}
{6}
{5}
{8}

\mcq
{116}
{দুটি Value সমান কি না যাচাই করতে কোন Operator ব্যবহৃত হয়?}
{=}
{==}
{!=}
{+=}

\mcq
{117}
{দুটি Value অসমান কি না যাচাই করতে কোন Operator ব্যবহৃত হয়?}
{==}
{!=}
{=}
{<=}

\mcq
{118}
{দুটি Condition-ই True হলে True পেতে কোন Logical Operator ব্যবহৃত হয়?}
{or}
{not}
{in}
{and}

\mcq
{119}
{একাধিক Condition-এর অন্তত একটি True হলেই True পেতে কোন Operator ব্যবহৃত হয়?}
{and}
{or}
{is}
{not}

\mcq
{120}
{Logical Value উল্টে দিতে কোন Operator ব্যবহৃত হয়?}
{not}
{or}
{and}
{+}

\mcqwide
{121}
{\textenglish{\ttfamily 2 + 3 * 4}-এর ফল কত?}
{24}
{14}
{9}
{20}

\mcqwide
{122}
{\textenglish{\ttfamily x += 2} কোনটির সমতুল্য?}
{x = 2}
{x = x + 2}
{x = x * 2}
{x == x + 2}

\topicheading{G. Conditional Logic Statement}
\mcq
{123}
{একটি Condition True হলে নির্দিষ্ট কাজ করার জন্য কোন Statement ব্যবহৃত হয়?}
{while}
{def}
{for}
{if}

\mcq
{124}
{Condition True ও False—দুই অবস্থায় আলাদা কাজ করতে কোন Structure উপযোগী?}
{import}
{def মাত্র}
{if-else}
{for মাত্র}

\mcq
{125}
{একাধিক Condition ধারাবাহিকভাবে যাচাই করতে কোনটি ব্যবহৃত হয়?}
{if-elif-else}
{while-pass}
{for-break}
{def-return}

\mcq
{126}
{Python-এ \py{if} Header-এর শেষে কোন Symbol থাকে?}
{;}
{:}
{,}
{.}

\mcq
{127}
{Conditional Block চিহ্নিত করতে Python-এ কোনটি গুরুত্বপূর্ণ?}
{Indentation}
{Curly brace বাধ্যতামূলক}
{Line number}
{Semicolon বাধ্যতামূলক}

\mcq
{128}
{কোন Condition দিয়ে Number জোড় কি না পরীক্ষা করা যায়?}
{n + 2 == 0}
{n == 2 সবসময়}
{n \char37 2 == 0}
{n / 2 == 1}

\mcq
{129}
{Positive Number শনাক্ত করতে কোন Condition সঠিক?}
{n > 0}
{n < 0}
{n != n}
{n == 0}

\mcq
{130}
{একটি \py{if} Block-এর ভিতরে আরেকটি \py{if} থাকলে তাকে কী বলে?}
{Global if}
{File if}
{Infinite if}
{Nested if}

\mcq
{131}
{একটি Number 10 থেকে 20-এর মধ্যে আছে কি না পরীক্ষা করতে কোনটি উপযোগী?}
{n == 10 and 20}
{not n}
{n >= 10 and n <= 20}
{n >= 10 or n <= 20}

\mcq
{132}
{Grade নির্ধারণের মতো একাধিক Range Condition-এর জন্য কোনটি বেশি উপযোগী?}
{if-elif-else}
{শুধু if একটি}
{শুধু break}
{শুধু pass}

\mcq
{133}
{Assignment ও Comparison-এর সঠিক জোড়া কোনটি?}
{= এবং ==}
{// এবং /}
{== এবং =}
{!= এবং +=}

\mcqwide
{134}
{\textenglish{\ttfamily if 0:} Condition-এর Boolean ফল কী?}
{Error সবসময়}
{True}
{None}
{False}

\mcqwide
{135}
{নিচের Code-এর Output কী?\codeblock{x = 7\\if x > 5:\\\hspace*{4mm}print("A")}}
{Error}
{5}
{কিছুই নয়}
{A}

\mcqwide
{136}
{Condition-এ \py{and} ব্যবহারের অর্থ কী?}
{দুই Condition-ই True হওয়া দরকার}
{Condition বাদ দেওয়া}
{সব Condition False হওয়া দরকার}
{একটি True হলেই যথেষ্ট}

\mcqwide
{137}
{Condition-এ \py{or} ব্যবহারের অর্থ কী?}
{সবসময় False}
{শুধু Number-এর জন্য}
{সব Condition True আবশ্যক}
{অন্তত একটি Condition True হলেই যথেষ্ট}

\mcqwide
{138}
{নিচের কোনটি Conditional Logic-এর সঠিক ব্যবহার?}
{একই Statement 100 বার চালানো}
{File Rename করা}
{শুধু Function Define করা}
{বয়স অনুযায়ী Category নির্ধারণ}

\topicheading{H. Loop, for, while, Loop Control ও Pattern}
\mcqwide
{139}
{Loop-এর প্রধান উদ্দেশ্য কী?}
{শুধু Function Define করা}
{Data Type পরিবর্তন}
{শুধু File খোলা}
{একই কাজ প্রয়োজনমতো বারবার করা}

\mcq
{140}
{\py{for} Loop সাধারণত কোনটির উপর Iteration করতে পারে?}
{শুধু File Name}
{Iterable}
{শুধু Boolean}
{শুধু Compiler}

\mcq
{141}
{\py{range(1, 5)}-এ Stop Value 5 কি অন্তর্ভুক্ত?}
{না}
{শুধু Even হলে}
{শুধু while-এ}
{হ্যাঁ}

\mcq
{142}
{\py{range(5)} থেকে কোন Sequence পাওয়া যায়?}
{5 মাত্র}
{0,1,2,3,4}
{0,1,2,3,4,5}
{1,2,3,4,5}

\mcq
{143}
{\py{range(1, 5)} থেকে কোন Sequence পাওয়া যায়?}
{1,2,3,4,5}
{1,3,5}
{1,2,3,4}
{0,1,2,3,4}

\mcq
{144}
{\py{range(1, 6, 2)} থেকে কোন Sequence পাওয়া যায়?}
{1,3,5}
{2,4,6}
{1,3,5,7}
{1,2,3,4,5}

\mcq
{145}
{\py{range(5, 0, -1)} থেকে প্রথম দুটি Value কোনটি?}
{0,1}
{4,3 দিয়ে শুরু হয় না}
{5,4}
{5,6}

\mcqwide
{146}
{\py{while} Loop কখন পর্যন্ত চলে?}
{Condition True থাকা পর্যন্ত}
{Condition False থাকা পর্যন্ত সবসময়}
{শুধু 10 বার}
{শুধু একবার}

\mcqwide
{147}
{\py{while} Loop-এ সাধারণত কোন তিনটি বিষয় গুরুত্বপূর্ণ?}
{Initialization, Condition, Update}
{Set, Dict, Tuple}
{Input, File, Compiler}
{Class, Object, Module}

\mcq
{148}
{\py{while} Loop-এ Update ভুলে গেলে কী হতে পারে?}
{Variable মুছে যায়}
{সবসময় Syntax Error}
{Infinite Loop}
{Compiler তৈরি হয়}

\mcqwide
{149}
{\py{break} কী করে?}
{কোনো কাজ করে না}
{চলমান Loop সম্পূর্ণ বন্ধ করে}
{Variable Global করে}
{বর্তমান Iteration শুধু Skip করে}

\mcqwide
{150}
{\py{continue} কী করে?}
{Program Exit করে}
{Function Return করে}
{পুরো Loop বন্ধ করে}
{বর্তমান Iteration Skip করে}

\mcqwide
{151}
{\py{pass} Statement-এর কাজ কী?}
{Value Return করা}
{Input নেওয়া}
{Loop বন্ধ করা}
{Placeholder হিসেবে কোনো কাজ না করা}

\mcqwide
{152}
{Loop-এর \py{else} Block সাধারণত কখন চলে?}
{শুধু continue হলে}
{সবসময় break হলে}
{Loop স্বাভাবিকভাবে শেষ হলে}
{শুধু Infinite Loop-এ}

\mcq
{153}
{একটি Loop-এর ভিতরে আরেকটি Loop থাকলে তাকে কী বলে?}
{Nested Loop}
{Binary Loop}
{File Loop}
{Global Loop}

\mcq
{154}
{Outer Loop 3 বার এবং Inner Loop প্রতি বার 2 বার চললে Inner Body মোট কতবার চলবে?}
{5}
{2}
{6}
{3}

\mcq
{155}
{Sum বের করার আগে \py{total} সাধারণত কত দিয়ে Initialize করা হয়?}
{0}
{-1}
{None সবসময়}
{1}

\mcq
{156}
{Factorial বের করার আগে \py{fact} সাধারণত কত দিয়ে Initialize করা হয়?}
{-1}
{0}
{2}
{1}

\mcqwide
{157}
{নিচের Code-এ \py{print()} মোট কয়বার Execute হবে?\codeblock{for i in range(1, 4):\\\hspace*{4mm}print(i)}}
{3}
{4}
{1}
{2}

\mcqwide
{158}
{নিচের Loop মোট কয়বার চলবে?\codeblock{i = 1\\while i <= 3:\\\hspace*{4mm}i += 1}}
{4}
{অসীম}
{2}
{3}

\mcq
{159}
{\py{continue} ব্যবহারের পর Loop কোথায় যায়?}
{Function-এর বাইরে}
{Loop-এর বাইরে}
{Program-এর শুরুতে}
{পরবর্তী Iteration-এ}

\mcqwide
{160}
{\py{break} Execute হওয়ার পর Control কোথায় যায়?}
{Loop-এর প্রথম Statement-এ}
{Current Iteration-এর শুরুতে}
{Program সবসময় বন্ধ হয়}
{Loop-এর পরের Statement-এ}

\mcq
{161}
{Pattern Printing-এ একই Line-এ একাধিক Item রাখতে \py{print()}-এর কোন Argument কাজে লাগে?}
{type}
{end}
{range}
{sep শুধু}

\mcq
{162}
{Pattern-এর একটি Row শেষ করে নতুন Line নিতে সাধারণত কী করা হয়?}
{আরেকটি print() ব্যবহার}
{break ব্যবহার}
{input() ব্যবহার}
{return ব্যবহার}

\mcqwide
{163}
{নিচের Code শেষে \py{total}-এর মান কত?\codeblock{total = 0\\for i in range(1, 5):\\\hspace*{4mm}total += i}}
{15}
{4}
{10}
{6}

\mcqwide
{164}
{\textenglish{\ttfamily for i in range(2,8,2)} Loop মোট কয়বার চলবে?}
{4}
{2}
{6}
{3}

\topicheading{I. Data Structure: List, Tuple, Set ও Dictionary}
\mcq
{165}
{List-এর বৈশিষ্ট্য কোনটি?}
{Unordered এবং শুধু unique}
{Ordered ও Mutable}
{শুধু key-value}
{Ordered ও Immutable}

\mcq
{166}
{Tuple-এর বৈশিষ্ট্য কোনটি?}
{শুধু Number}
{Unique এবং indexing নেই}
{Ordered ও Immutable}
{Mutable ও key-value}

\mcqwide
{167}
{Set-এর প্রধান বৈশিষ্ট্য কোনটি?}
{Unique Element}
{Duplicate বাধ্যতামূলক}
{Key-value pair}
{Indexed sequence হিসেবে ব্যবহার বাধ্যতামূলক}

\mcq
{168}
{Dictionary কীভাবে Data সংরক্ষণ করে?}
{শুধু Tuple}
{শুধু Index}
{শুধু Value}
{Key–Value Pair}

\mcq
{169}
{List-এর প্রথম Element-এর Index কত?}
{1}
{2}
{-1}
{0}

\mcq
{170}
{List বা Tuple-এর শেষ Element পেতে কোন Negative Index ব্যবহৃত হয়?}
{1}
{0}
{-2}
{-1}

\mcq
{171}
{List-এর শেষে নতুন Element যোগ করতে কোন Method ব্যবহৃত হয়?}
{add()}
{push()}
{append()}
{join()}

\mcq
{172}
{List-এর নির্দিষ্ট Position-এ Element যোগ করতে কোন Method ব্যবহৃত হয়?}
{range()}
{type()}
{insert()}
{append()}

\mcq
{173}
{List থেকে নির্দিষ্ট Value সরাতে কোন Method ব্যবহার করা যায়?}
{insert()}
{append()}
{type()}
{remove()}

\mcq
{174}
{Single-item Tuple-এর সঠিক রূপ কোনটি?}
{(10)}
{[10]}
{(10,)}
{{10:}}

\mcq
{175}
{Empty Set তৈরির সঠিক উপায় কোনটি?}
{[]}
{()}
{{}}
{set()}

\mcqwide
{176}
{Python-এ খালি \textenglish{\ttfamily \{\}} কী তৈরি করে?}
{List}
{Dictionary}
{Set}
{Tuple}

\mcq
{177}
{Dictionary-তে \py{in} Operator সাধারণত কী পরীক্ষা করে?}
{Index}
{Value মাত্র}
{Key}
{Tuple Length}

\mcq
{178}
{Collection-এর Element সংখ্যা জানার জন্য কোন Function ব্যবহৃত হয়?}
{input()}
{len()}
{type()}
{range()}

\mcq
{179}
{List-এর একটি অংশ নেওয়ার কৌশলকে কী বলা হয়?}
{Casting}
{Tracing}
{Compiling}
{Slicing}

\mcq
{180}
{নিচের কোনটি Mutable?}
{Integer}
{Tuple}
{List}
{String}

\mcq
{181}
{নিচের কোনটি Immutable?}
{Tuple}
{Set}
{Dictionary}
{List}

\mcq
{182}
{Duplicate বাদ দিয়ে Unique Value রাখতে কোনটি বেশি উপযোগী?}
{List}
{Tuple}
{Set}
{String}

\topicheading{J. Function, Function Contract, Scope, Local/Global ও File}
\mcqwide
{183}
{Function কী?}
{শুধু Loop}
{নির্দিষ্ট কাজের reusable Code Block}
{শুধু Variable}
{শুধু Data Type}

\mcq
{184}
{Python-এ Function Define করতে কোন Keyword ব্যবহৃত হয়?}
{def}
{function}
{return}
{fun}

\mcq
{185}
{Function ব্যবহার করতে তাকে কী করতে হয়?}
{Global করতে}
{Compile আলাদা করে}
{Call}
{সবসময় File-এ লিখতে}

\mcq
{186}
{Function Definition-এ Input গ্রহণকারী Variable-কে কী বলে?}
{Argument}
{Scope}
{Parameter}
{Return}

\mcq
{187}
{Function Call-এর সময় পাঠানো Actual Value-কে কী বলে?}
{Parameter}
{Argument}
{Local}
{Contract}

\mcqwide
{188}
{\py{return}-এর প্রধান কাজ কী?}
{Variable Delete করা}
{Function থেকে Value ফেরত পাঠানো}
{Loop বন্ধ করা}
{শুধু Screen-এ Print}

\mcqwide
{189}
{\py{print()} ও \py{return}-এর মূল পার্থক্য কোনটি?}
{print Function বন্ধ করে}
{print Output দেখায়, return Value ফেরত দেয়}
{return শুধু Comment}
{দুটিই সবসময় একই}

\mcq
{190}
{Parameter-এর Default Value থাকলে তাকে কী বলা হয়?}
{Loop Parameter}
{File Parameter}
{Global Parameter}
{Default Parameter}

\mcq
{191}
{Function Call-এ Parameter Name উল্লেখ করে Value পাঠানোকে কী বলা হয়?}
{Keyword Argument}
{Machine Argument}
{Binary Argument}
{Loop Argument}

\mcq
{192}
{কোনো Function-এ \py{return} না থাকলে সাধারণত কী Return হয়?}
{Error সবসময়}
{0 সবসময়}
{None}
{False সবসময়}

\mcq
{193}
{Function-এর ভিতরে তৈরি Variable সাধারণত কোন Scope-এর?}
{Universal}
{Global}
{Local}
{File}

\mcq
{194}
{Function-এর বাইরে Module Level-এ তৈরি Variable সাধারণত কী?}
{Local Variable}
{Global Variable}
{Argument}
{Parameter}

\mcq
{195}
{Function-এর ভিতরে Global Variable-এ Assignment করতে কোন Keyword দরকার হতে পারে?}
{global}
{pass}
{break}
{return}

\mcqwide
{196}
{Scope কী নির্ধারণ করে?}
{কোথায় একটি Name/Variable ব্যবহারযোগ্য}
{Variable-এর Font}
{File Extension}
{CPU Speed}

\mcq
{197}
{Function-এর Local Variable-এর নাম Global Variable-এর মতো হলে সাধারণত Function-এর ভিতরে কোনটি ব্যবহৃত হয়?}
{দুটির যোগফল}
{Local Variable}
{Global Variable বাধ্যতামূলক}
{কোনোটিই নয়}

\mcqwide
{198}
{Function Contract মূলত কী বোঝায়?}
{Function কী Input নেবে, কী করবে ও কী Result দেবে}
{শুধু Function-এর নাম}
{শুধু Loop count}
{শুধু File Name}

\mcq
{199}
{Program বন্ধ হওয়ার পরও Data রাখতে কোনটি প্রয়োজন হতে পারে?}
{File}
{Parameter}
{Local Variable মাত্র}
{Loop Counter}

\mcqwide
{200}
{Variable ও File-এর মূল পার্থক্য কোনটি?}
{File শুধু Loop}
{Variable Runtime Data রাখে, File persistent Data রাখতে পারে}
{Variable সবসময় Disk-এ থাকে}
{দুটিই একই}

\mcqwide
{201}
{অপ্রয়োজনীয় Global Variable কম ব্যবহার করার একটি কারণ কী?}
{সব Function বন্ধ করা}
{Dependency ও Side Effect কমানো}
{Loop দ্রুত করা সবসময়}
{Compiler বাদ দেওয়া}

\mcqwide
{202}
{নিচের কোন বক্তব্যটি সঠিক?\newline i. Parameter Function Definition-এ থাকে।\newline ii. Argument Function Call-এ পাঠানো হয়।\newline iii. return Result ফেরত দিতে পারে।}
{i ও iii}
{ii ও iii}
{i ও ii}
{i, ii ও iii}

\topicheading{N. Python Code Practice — Output, Dry Run ও Trap}

\mcqwide
{203}
{নিচের Code-এর Output কত?\codeblock{x = 5\\y = 7\\print(x + y)}}
{10}
{12}
{57}
{Error}

\mcqwide
{204}
{নিচের Code-এর Output কোনটি?\codeblock{x = "5"\\y = "2"\\print(x + y)}}
{7}
{52}
{10}
{Error}

\mcqwide
{205}
{নিচের Code-এর Output কত?\codeblock{print(2 + 3 * 4)}}
{20}
{14}
{24}
{10}

\mcqwide
{206}
{নিচের Code-এর Output কী?\codeblock{print(not (5 > 3))}}
{True}
{False}
{5}
{Error}

\mcqwide
{207}
{নিচের Code-এর Output কোনটি?\codeblock{x = 0\\if x:\\\hspace*{4mm}print("Yes")\\else:\\\hspace*{4mm}print("No")}}
{Yes}
{No}
{0}
{Error}

\mcqwide
{208}
{নিচের Code-এর Output কোনটি?\codeblock{for i in range(5, 0, -2):\\\hspace*{4mm}print(i, end=" ")}}
{5 3 1}
{5 4 3 2 1}
{1 3 5}
{5 3}

\mcqwide
{209}
{নিচের Code শেষে \py{total}-এর মান কত?\codeblock{total = 0\\for i in range(2, 7, 2):\\\hspace*{4mm}total += i\\print(total)}}
{6}
{8}
{12}
{20}

\mcqwide
{210}
{নিচের Code-এর Output কোনটি?\codeblock{i = 1\\while i <= 4:\\\hspace*{4mm}print(i, end=" ")\\\hspace*{4mm}i += 2}}
{1 2 3 4}
{1 3}
{1 3 5}
{2 4}

\mcqwide
{211}
{নিচের Code-এ \py{print()} কতবার Execute হবে?\codeblock{for i in range(1, 6):\\\hspace*{4mm}if i == 3:\\\hspace*{8mm}continue\\\hspace*{4mm}print(i)}}
{3 বার}
{4 বার}
{5 বার}
{2 বার}

\mcqwide
{212}
{নিচের Code-এ \py{print("*")} মোট কতবার Execute হবে?\codeblock{for i in range(2):\\\hspace*{4mm}for j in range(3):\\\hspace*{8mm}print("*")}}
{2 বার}
{3 বার}
{5 বার}
{6 বার}

\mcqwide
{213}
{নিচের Code-এর Output কত?\codeblock{a = [10, 20, 30, 40]\\print(a[-1])}}
{10}
{20}
{30}
{40}

\mcqwide
{214}
{নিচের Code-এর পরে List কোনটি হবে?\codeblock{a = [1, 2]\\a.append(3)\\a.pop(0)\\print(a)}}
{[1, 2, 3]}
{[2, 3]}
{[1, 3]}
{[2]}

\mcqwide
{215}
{নিচের Code-এর Output কত?\codeblock{s = \{1, 2, 2, 3\}\\print(len(s))}}
{2}
{3}
{4}
{Error}

\mcqwide
{216}
{নিচের Code-এর Output কোনটি?\codeblock{d = \{"name": "Rafi", "mark": 80\}\\print(d["mark"])}}
{Rafi}
{mark}
{80}
{Error}

\mcqwide
{217}
{নিচের Code-এর Output কত?\codeblock{def add(a, b=2):\\\hspace*{4mm}return a + b\\print(add(5))}}
{2}
{5}
{7}
{Error}

\mcqwide
{218}
{নিচের Code-এর Output কত?\codeblock{def square(x):\\\hspace*{4mm}return x * x\\y = square(4) + 1\\print(y)}}
{16}
{17}
{8}
{5}

\mcqwide
{219}
{নিচের Code-এর Output কোনটি?\codeblock{x = 10\\def show():\\\hspace*{4mm}x = 20\\\hspace*{4mm}print(x)\\show()\\print(x)}}
{10 10}
{20 20}
{20 10}
{Error}

\mcqwide
{220}
{নিচের Code-এর Output কত?\codeblock{def double(x):\\\hspace*{4mm}return x * 2\\def calc(n):\\\hspace*{4mm}return double(n) + 1\\print(calc(3))}}
{6}
{7}
{8}
{Error}

\end{multicols}
\topicheading{Topic-wise MCQ Answer Key}
\begingroup
\setlength{\tabcolsep}{2.4pt}
\renewcommand{\arraystretch}{1.08}
\begin{longtable}{|C{0.048\linewidth}|C{0.037\linewidth}|C{0.048\linewidth}|C{0.037\linewidth}|C{0.048\linewidth}|C{0.037\linewidth}|C{0.048\linewidth}|C{0.037\linewidth}|C{0.048\linewidth}|C{0.037\linewidth}|C{0.048\linewidth}|C{0.037\linewidth}|C{0.048\linewidth}|C{0.037\linewidth}|}
\hline
\textbf{নং} & \textbf{উঃ} & \textbf{নং} & \textbf{উঃ} & \textbf{নং} & \textbf{উঃ} & \textbf{নং} & \textbf{উঃ} & \textbf{নং} & \textbf{উঃ} & \textbf{নং} & \textbf{উঃ} & \textbf{নং} & \textbf{উঃ} \\
\hline
\endfirsthead
\hline
\rowcolor{softBlue}
\textbf{নং} & \textbf{উঃ} & \textbf{নং} & \textbf{উঃ} & \textbf{নং} & \textbf{উঃ} & \textbf{নং} & \textbf{উঃ} & \textbf{নং} & \textbf{উঃ} & \textbf{নং} & \textbf{উঃ} & \textbf{নং} & \textbf{উঃ} \\
\hline
\endhead
1 & গ & 2 & ঘ & 3 & গ & 4 & ঘ & 5 & খ & 6 & ঘ & 7 & ঘ \\
\hline
8 & খ & 9 & খ & 10 & গ & 11 & ক & 12 & গ & 13 & ক & 14 & গ \\
\hline
15 & ক & 16 & খ & 17 & ক & 18 & খ & 19 & ঘ & 20 & ক & 21 & খ \\
\hline
22 & গ & 23 & ক & 24 & ঘ & 25 & ঘ & 26 & ঘ & 27 & খ & 28 & ঘ \\
\hline
29 & ক & 30 & ক & 31 & গ & 32 & গ & 33 & ক & 34 & খ & 35 & ক \\
\hline
36 & ঘ & 37 & ঘ & 38 & গ & 39 & ঘ & 40 & গ & 41 & ক & 42 & গ \\
\hline
43 & ঘ & 44 & ক & 45 & খ & 46 & গ & 47 & ঘ & 48 & গ & 49 & খ \\
\hline
50 & গ & 51 & ক & 52 & খ & 53 & ক & 54 & ঘ & 55 & গ & 56 & গ \\
\hline
57 & খ & 58 & গ & 59 & ঘ & 60 & ক & 61 & ক & 62 & গ & 63 & গ \\
\hline
64 & খ & 65 & ঘ & 66 & খ & 67 & গ & 68 & ক & 69 & খ & 70 & ক \\
\hline
71 & খ & 72 & খ & 73 & খ & 74 & খ & 75 & গ & 76 & খ & 77 & ঘ \\
\hline
78 & গ & 79 & ক & 80 & ক & 81 & ক & 82 & গ & 83 & গ & 84 & ক \\
\hline
85 & গ & 86 & গ & 87 & ঘ & 88 & খ & 89 & ঘ & 90 & ঘ & 91 & ক \\
\hline
92 & গ & 93 & ক & 94 & ক & 95 & খ & 96 & ঘ & 97 & ঘ & 98 & ঘ \\
\hline
99 & গ & 100 & খ & 101 & গ & 102 & খ & 103 & ক & 104 & ক & 105 & খ \\
\hline
106 & খ & 107 & ঘ & 108 & গ & 109 & গ & 110 & ঘ & 111 & ঘ & 112 & খ \\
\hline
113 & ক & 114 & খ & 115 & ঘ & 116 & খ & 117 & খ & 118 & ঘ & 119 & খ \\
\hline
120 & ক & 121 & খ & 122 & খ & 123 & ঘ & 124 & গ & 125 & ক & 126 & খ \\
\hline
127 & ক & 128 & গ & 129 & ক & 130 & ঘ & 131 & গ & 132 & ক & 133 & ক \\
\hline
134 & ঘ & 135 & ঘ & 136 & ক & 137 & ঘ & 138 & ঘ & 139 & ঘ & 140 & খ \\
\hline
141 & ক & 142 & খ & 143 & গ & 144 & ক & 145 & গ & 146 & ক & 147 & ক \\
\hline
148 & গ & 149 & খ & 150 & ঘ & 151 & ঘ & 152 & গ & 153 & ক & 154 & গ \\
\hline
155 & ক & 156 & ঘ & 157 & ক & 158 & ঘ & 159 & ঘ & 160 & ঘ & 161 & খ \\
\hline
162 & ক & 163 & গ & 164 & ঘ & 165 & খ & 166 & গ & 167 & ক & 168 & ঘ \\
\hline
169 & ঘ & 170 & ঘ & 171 & গ & 172 & গ & 173 & ঘ & 174 & গ & 175 & ঘ \\
\hline
176 & খ & 177 & গ & 178 & খ & 179 & ঘ & 180 & গ & 181 & ক & 182 & গ \\
\hline
183 & খ & 184 & ক & 185 & গ & 186 & গ & 187 & খ & 188 & খ & 189 & খ \\
\hline
190 & ঘ & 191 & ক & 192 & গ & 193 & গ & 194 & খ & 195 & ক & 196 & ক \\
\hline
197 & খ & 198 & ক & 199 & ক & 200 & খ & 201 & খ & 202 & ঘ & 203 & খ \\
\hline
204 & খ & 205 & খ & 206 & খ & 207 & খ & 208 & ক & 209 & গ & 210 & খ \\
\hline
211 & খ & 212 & ঘ & 213 & ঘ & 214 & খ & 215 & খ & 216 & গ & 217 & গ \\
\hline
218 & খ & 219 & গ & 220 & খ & -- & -- & -- & -- & -- & -- & -- & -- \\
\hline
\end{longtable}
\endgroup
\partheading{অংশ–২ : Board-style Mixed MCQ}
\begin{infobox}
এখানে Topic-এর নাম দেওয়া নেই। Direct Concept, Statement, Python Code, Output, Dry Run, Trap ও Application-based প্রশ্ন মিশিয়ে রাখা হয়েছে, যাতে পরীক্ষার মতো নিজেই প্রয়োজনীয় Concept শনাক্ত করতে হয়।
\end{infobox}
\begin{multicols}{2}
\setlength{\columnsep}{4mm}
\raggedcolumns
\mcqwide
{1}
{Machine Language থেকে High Level Language-এ যাওয়ার প্রধান সুবিধা কোনটি সবচেয়ে যথাযথ?}
{Program লেখা ও বোঝা সহজ হয়}
{কম্পিউটার Binary বোঝে না}
{Hardware dependence বাড়ে}
{Assembler প্রয়োজন হয় বেশি}

\mcqwide
{2}
{নিচের বক্তব্যগুলো লক্ষ করো—\newline i. Machine Language machine-dependent।\newline ii. Assembly Language-এ mnemonic থাকে।\newline iii. High Level Language সাধারণত মানুষের জন্য সহজ।\newline সঠিক কোনটি?}
{i ও iii}
{i, ii ও iii}
{i ও ii}
{ii ও iii}

\mcq
{3}
{একজন Programmer এমন Source Code লিখল যা পুরোটা একসঙ্গে অনুবাদ হওয়ার পর Execute হবে। এখানে Translator কোনটি?}
{Interpreter}
{Assembler}
{Compiler}
{Editor}

\mcq
{4}
{একটি Program Statement ধরে ধরে Execute হচ্ছে এবং Error পাওয়া মাত্র থেমে যাচ্ছে। কোন Translator-এর বৈশিষ্ট্য?}
{Compiler}
{Assembler}
{Interpreter}
{Loader}

\mcq
{5}
{Assembly Code-কে সরাসরি Machine Code-এ রূপান্তরের জন্য কোনটি প্রয়োজন?}
{Interpreter}
{Assembler}
{Compiler}
{Debugger}

\mcq
{6}
{একটি সমস্যায় Input হিসেবে দৈর্ঘ্য ও প্রস্থ নেওয়া হলো, তারপর গুণ করে ক্ষেত্রফল দেখানো হলো। IPO Model-এ “গুণ” কোন অংশ?}
{Storage}
{Processing}
{Output}
{Input}

\mcq
{7}
{একজন শিক্ষার্থী Program না চালিয়ে Variable-এর মান প্রতি ধাপে লিখে Result যাচাই করছে। সে কোন কৌশল ব্যবহার করছে?}
{Installation}
{Documentation}
{Compilation}
{Dry Run}

\mcq
{8}
{Flowchart-এ একটি Number জোড় কি না যাচাইয়ের স্থানে কোন Symbol সবচেয়ে উপযুক্ত?}
{Diamond}
{Parallelogram}
{Rectangle}
{Oval}

\mcq
{9}
{নিচের কোন Flowchart mapping সঠিক?}
{Input—Parallelogram}
{Process—Oval}
{Start—Diamond}
{Decision—Rectangle}

\mcq
{10}
{Algorithm-এর কোন বৈশিষ্ট্য না থাকলে তা কার্যকর Algorithm হিসেবে গ্রহণযোগ্য নয়?}
{সসীম ধাপে শেষ হওয়া}
{Python Syntax থাকা}
{শুধু English-এ লেখা}
{শুধু Diagram হওয়া}

\mcq
{11}
{নিচের কোন Variable Name-টি Python-এ বৈধ?}
{for}
{2result}
{final-result}
{final\_result}

\mcq
{12}
{নিচের কোন জোড়াটি একই Variable নির্দেশ করে না?}
{name ও name}
{x ও x}
{total ও total}
{score ও Score}

\mcqwide
{13}
{\textenglish{\ttfamily value = input("Enter: ")} দ্বারা পাওয়া \py{value}-এর Type সাধারণত কী?}
{str}
{bool}
{int}
{float}

\mcq
{14}
{কোন Statement-এ Input নেওয়ার সাথে সাথে Integer-এ Convert করা হয়েছে?}
{\textenglish{\ttfamily n = type(input())}}
{\textenglish{\ttfamily n = input()}}
{\textenglish{\ttfamily n = int(input())}}
{\textenglish{\ttfamily n = print()}}

\mcqwide
{15}
{\textenglish{\ttfamily x = 5 + 2 * 3} হলে x-এর মান কত?}
{15}
{13}
{11}
{21}

\mcqwide
{16}
{নিচের Code-এ \py{a+b}-এর মান কত?\codeblock{a = 17 // 5\\b = 17 \char37 5}}
{5}
{4}
{7}
{3}

\mcqwide
{17}
{\textenglish{\ttfamily 3 ** 2 + 1}-এর ফল কত?}
{6}
{9}
{7}
{10}

\mcqwide
{18}
{নিচের বক্তব্যগুলো লক্ষ করো—\newline i. \py{=} Assignment Operator।\newline ii. \py{==} Comparison Operator।\newline iii. \py{!=} অসমতা পরীক্ষা করে।\newline সঠিক কোনটি?}
{i ও ii}
{i ও iii}
{i, ii ও iii}
{ii ও iii}

\mcqwide
{19}
{নিচের Code-এর Output কী?\codeblock{x = 8\\if x \char37 2 == 0:\\\hspace*{4mm}print("Even")}}
{8}
{Odd}
{Error}
{Even}

\mcqwide
{20}
{নিচের Code-এর Output কী?\codeblock{x = 5\\if x > 5:\\\hspace*{4mm}print("A")\\else:\\\hspace*{4mm}print("B")}}
{B}
{কিছুই নয়}
{5}
{A}

\mcq
{21}
{একটি Mark 80 বা তার বেশি হলে A+, 70 বা তার বেশি হলে A—এভাবে একাধিক Range যাচাইয়ে কোন Structure সবচেয়ে উপযুক্ত?}
{শুধু for}
{শুধু pass}
{if-elif-else}
{শুধু while}

\mcqwide
{22}
{নিচের Code-এর Output কী?\codeblock{n = 15\\if n >= 10 and n <= 20:\\\hspace*{4mm}print("Yes")}}
{Error}
{Yes}
{15}
{No}

\mcqwide
{23}
{নিচের Code-এ মোট কয়টি Number Print হবে?\codeblock{for i in range(1, 6):\\\hspace*{4mm}print(i, end=" ")}}
{1}
{5}
{4}
{6}

\mcqwide
{24}
{\textenglish{\ttfamily range(2, 10, 2)} থেকে কোন Sequence পাওয়া যায়?}
{2,3,4,5,6,7,8,9}
{0,2,4,6,8}
{2,4,6,8,10}
{2,4,6,8}

\mcqwide
{25}
{\textenglish{\ttfamily range(5, 0, -2)} থেকে কোন Sequence পাওয়া যায়?}
{0,2,4}
{5,3,1}
{5,3,1,-1}
{5,4,3,2,1}

\mcqwide
{26}
{নিচের Code শেষে \py{total}-এর মান কত?\codeblock{total = 0\\for i in range(1, 5):\\\hspace*{4mm}total += i}}
{6}
{10}
{15}
{5}

\mcqwide
{27}
{নিচের Code শেষে \py{fact}-এর মান কত?\codeblock{fact = 1\\for i in range(1, 5):\\\hspace*{4mm}fact *= i}}
{120}
{24}
{10}
{4}

\mcqwide
{28}
{নিচের Code-এর Output কোনটি?\codeblock{for i in range(1, 6):\\\hspace*{4mm}if i == 3:\\\hspace*{8mm}break\\\hspace*{4mm}print(i, end=" ")}}
{1 2}
{3 4 5}
{1 2 4 5}
{1 2 3}

\mcqwide
{29}
{নিচের Code-এর Output কোনটি?\codeblock{for i in range(1, 6):\\\hspace*{4mm}if i == 3:\\\hspace*{8mm}continue\\\hspace*{4mm}print(i, end=" ")}}
{1 2}
{1 2 4 5}
{1 2 3 4 5}
{3}

\mcqwide
{30}
{নিচের Code-এর Output কী?\codeblock{i = 1\\while i <= 3:\\\hspace*{4mm}print(i, end=" ")\\\hspace*{4mm}i += 1}}
{Infinite}
{1 2 3 4}
{1 2 3}
{0 1 2 3}

\mcq
{31}
{একটি \py{while} Loop-এ Counter Update নেই এবং Condition সবসময় True থাকে। সম্ভাব্য ফল কী?}
{Infinite Loop}
{Function Return করে}
{Loop একবার চলে}
{Syntax সবসময় ঠিক হয় না}

\mcq
{32}
{Outer Loop 4 বার ও Inner Loop 3 বার চললে Inner Body মোট কতবার Execute হবে?}
{7}
{12}
{3}
{4}

\mcq
{33}
{Pattern Printing-এ Inner Loop-এর \py{print()}-এ \py{end=" "} ব্যবহারের উদ্দেশ্য কী?}
{Input নেওয়া}
{Function Return করা}
{একই Row-তে Output রাখা}
{Program বন্ধ করা}

\mcq
{34}
{নিচের কোন Collection-টি পরিবর্তনযোগ্য এবং Index ব্যবহার করে Element Access করা যায়?}
{Dictionary শুধু}
{Set}
{List}
{Tuple}

\mcq
{35}
{একটি Collection-এ Duplicate বাদ দিয়ে শুধু Unique Roll Number রাখতে চাই। কোনটি উপযুক্ত?}
{Set}
{List}
{Tuple}
{String}

\mcq
{36}
{একটি Student-এর Name, Roll, GPA-কে label সহ সংরক্ষণে কোন Data Structure সবচেয়ে উপযোগী?}
{Dictionary}
{String}
{Tuple মাত্র}
{Set}

\mcqwide
{37}
{\textenglish{\ttfamily a = [10,20,30]} হলে \textenglish{\ttfamily a[-1]} কত?}
{20}
{10}
{30}
{Error}

\mcqwide
{38}
{নিচের Code-এর পরে Final List কোনটি?\codeblock{a = [1, 2]\\a.append(3)}}
{[1,2,3]}
{[1,3]}
{[1,2]}
{[3,1,2]}

\mcq
{39}
{নিচের কোনটি Single-item Tuple?}
{set(5)}
{(5)}
{[5]}
{(5,)}

\mcq
{40}
{খালি Set প্রয়োজন হলে কোনটি সঠিক?}
{()}
{[]}
{set()}
{{}}

\mcqwide
{41}
{নিচের বক্তব্যগুলো লক্ষ করো—\newline i. List Mutable।\newline ii. Tuple Immutable।\newline iii. Set Unique Element রাখে।\newline সঠিক কোনটি?}
{i ও ii}
{ii ও iii}
{i, ii ও iii}
{i ও iii}

\mcqwide
{42}
{একটি Dictionary \textenglish{\ttfamily d = \{"a":1, "b":2\}} হলে \textenglish{\ttfamily "a" in d} এর ফল কী?}
{Error}
{False}
{True}
{1}

\mcqwide
{43}
{নিচের Function-এ \py{a} ও \py{b} কী?\codeblock{def add(a, b):\\\hspace*{4mm}return a + b}}
{Parameter}
{File}
{Argument}
{Global Variable}

\mcqwide
{44}
{\textenglish{\ttfamily add(2,3)} Call-এ 2 ও 3 কী?}
{Keyword}
{Scope}
{Parameter}
{Argument}

\mcqwide
{45}
{নিচের Code-এ \py{y}-এর মান কত?\codeblock{def f(x):\\\hspace*{4mm}return x * 2\\y = f(4)}}
{4}
{2}
{None}
{8}

\mcqwide
{46}
{নিচের Code-এ \py{x}-এর Value কী?\codeblock{def f():\\\hspace*{4mm}print("Hi")\\x = f()}}
{0}
{None}
{Hi}
{Error সবসময়}

\mcqwide
{47}
{নিচের Code-এর Output কোনটি?\codeblock{x = 10\\def show():\\\hspace*{4mm}x = 20\\\hspace*{4mm}print(x)\\show()\\print(x)}}
{Error}
{20 তারপর 10}
{10 তারপর 20}
{20 তারপর 20}

\mcq
{48}
{Function-এর ভিতরে বাইরের Global Variable-এর Value পরিবর্তন করতে কোন Keyword প্রয়োজন হতে পারে?}
{continue}
{pass}
{return}
{global}

\mcqwide
{49}
{Function Contract-এর কোনটি অংশ নয়?}
{Function-এর কাজ}
{Function-এর Input প্রত্যাশা}
{Expected Return/Output}
{Computer-এর Wallpaper}

\mcq
{50}
{Program বন্ধ হওয়ার পরও Result সংরক্ষণ করতে চাই। কোনটি সবচেয়ে উপযুক্ত?}
{Loop Counter}
{File}
{Parameter}
{Local Variable}

\mcq
{51}
{একটি Program-এ একই Calculation বারবার দরকার। কোন কৌশল Code Repetition কমাতে সবচেয়ে কার্যকর?}
{Global Variable বাড়ানো}
{শুধু Comment}
{Function}
{সব Code copy করা}

\mcqwide
{52}
{নিচের কোন পরিস্থিতিতে \py{continue} বেশি উপযোগী?}
{Loop পুরো বন্ধ করতে}
{File save করতে}
{কিছু নির্দিষ্ট Iteration বাদ দিয়ে বাকিগুলো চালাতে}
{Function Value ফেরাতে}

\mcqwide
{53}
{নিচের কোন পরিস্থিতিতে \py{break} বেশি উপযোগী?}
{Placeholder রাখতে}
{চাওয়া Value পাওয়া গেলে Search Loop থামাতে}
{Function Define করতে}
{একটি Iteration বাদ দিতে}

\mcq
{54}
{একটি Number 3 ও 5 উভয় দ্বারা বিভাজ্য কি না পরীক্ষা করতে কোন Condition সঠিক?}
{n \char37 3 == 0 or n \char37 5 == 0}
{n == 3 and 5}
{n / 3 == 5}
{n \char37 3 == 0 and n \char37 5 == 0}

\mcqwide
{55}
{নিচের Loop শেষে \py{x}-এর মান কত?\codeblock{x = 1\\while x < 10:\\\hspace*{4mm}x *= 2}}
{4}
{8}
{10}
{16}

\mcqwide
{56}
{\textenglish{\ttfamily nums = [2,4,6,8]} হলে \textenglish{\ttfamily len(nums)} কত?}
{8}
{2}
{4}
{3}

\mcqwide
{57}
{\textenglish{\ttfamily nums = [1,2,3,4]} হলে \textenglish{\ttfamily nums[1:3]} কোনটি?}
{[2,3,4]}
{[1,2]}
{[2,3]}
{[1,2,3]}

\mcq
{58}
{একটি List-এ Value আছে কি না দ্রুত যাচাইয়ের সহজ Expression কোনটি?}
{range(x)}
{x in mylist}
{x == mylist সবসময়}
{type(x)}

\mcq
{59}
{একটি Function-এর ভিতরে Global Variable পড়া হচ্ছে কিন্তু Assignment করা হচ্ছে না। সাধারণত \py{global} Keyword কি আবশ্যক?}
{হ্যাঁ, সবসময়}
{না}
{শুধু int হলে}
{শুধু return থাকলে}

\mcqwide
{60}
{নিচের কোনটি File ব্যবহারের সবচেয়ে যৌক্তিক কারণ?}
{Syntax Error দূর করা}
{Program শেষ হওয়ার পরও Data সংরক্ষণ}
{Loop-এর Iteration বাড়ানো}
{Variable-এর নাম ছোট করা}

\mcqwide
{61}
{নিচের Code চালালে কী হবে?\codeblock{x = "10"\\y = 2\\print(x + y)}}
{12 Print হবে}
{102 Print হবে}
{TypeError হবে}
{10 Print হবে}

\mcqwide
{62}
{নিচের Code-এর Output কী?\codeblock{x = 12\\if x \char37 2 == 0 and x > 10:\\\hspace*{4mm}print("A")\\else:\\\hspace*{4mm}print("B")}}
{A}
{B}
{12}
{Error}

\mcqwide
{63}
{নিচের Code-এর Output কোনটি?\codeblock{for i in range(6, 1, -2):\\\hspace*{4mm}print(i, end=" ")}}
{6 4 2}
{6 4}
{2 4 6}
{6 5 4 3 2}

\mcqwide
{64}
{নিচের Code-এর প্রধান সমস্যা কী?\codeblock{i = 1\\while i <= 5:\\\hspace*{4mm}print(i)}}
{Syntax Error}
{Loop একবার চলবে}
{Infinite Loop হতে পারে}
{কোনো Output হবে না}

\mcqwide
{65}
{নিচের Code-এ \py{count += 1} মোট কতবার Execute হবে?\codeblock{count = 0\\for i in range(3):\\\hspace*{4mm}for j in range(2):\\\hspace*{8mm}count += 1\\print(count)}}
{3 বার}
{5 বার}
{6 বার}
{9 বার}

\mcqwide
{66}
{নিচের Code-এর Output কোনটি?\codeblock{a = [10, 20, 30, 40, 50]\\print(a[1:4])}}
{[10, 20, 30]}
{[20, 30, 40]}
{[20, 30, 40, 50]}
{[10, 20, 30, 40]}

\mcqwide
{67}
{নিচের Code-এর Output কত?\codeblock{a = [1, 1, 2, 3, 3]\\s = set(a)\\print(len(s))}}
{2}
{3}
{5}
{Error}

\mcqwide
{68}
{নিচের Code-এর Output কী?\codeblock{d = \{"a": 10, "b": 20\}\\print("a" in d)}}
{True}
{False}
{10}
{Error}

\mcqwide
{69}
{নিচের Function Call-এ 5 ও 3 কী?\codeblock{def add(a, b):\\\hspace*{4mm}return a + b\\add(5, 3)}}
{Parameter}
{Argument}
{Local Function}
{Operator}

\mcqwide
{70}
{নিচের Code-এর Output কত?\codeblock{def power(x, p=2):\\\hspace*{4mm}return x ** p\\print(power(3))}}
{6}
{8}
{9}
{Error}

\mcqwide
{71}
{নিচের Code-এর Output কোনটি?\codeblock{x = 5\\def change():\\\hspace*{4mm}global x\\\hspace*{4mm}x = 9\\change()\\print(x)}}
{5}
{9}
{14}
{Error}

\mcqwide
{72}
{নিচের Code-এর Output-এর প্রকৃতি কোনটি?\codeblock{def show():\\\hspace*{4mm}print("Hi")\\x = show()\\print(x)}}
{শুধু Hi}
{Hi তারপর None}
{শুধু None}
{Error}

\end{multicols}
\topicheading{Board-style MCQ Answer Key}
\begingroup
\setlength{\tabcolsep}{2.4pt}
\renewcommand{\arraystretch}{1.08}
\begin{longtable}{|C{0.048\linewidth}|C{0.037\linewidth}|C{0.048\linewidth}|C{0.037\linewidth}|C{0.048\linewidth}|C{0.037\linewidth}|C{0.048\linewidth}|C{0.037\linewidth}|C{0.048\linewidth}|C{0.037\linewidth}|C{0.048\linewidth}|C{0.037\linewidth}|C{0.048\linewidth}|C{0.037\linewidth}|}
\hline
\textbf{নং} & \textbf{উঃ} & \textbf{নং} & \textbf{উঃ} & \textbf{নং} & \textbf{উঃ} & \textbf{নং} & \textbf{উঃ} & \textbf{নং} & \textbf{উঃ} & \textbf{নং} & \textbf{উঃ} & \textbf{নং} & \textbf{উঃ} \\
\hline
\endfirsthead
\hline
\rowcolor{softBlue}
\textbf{নং} & \textbf{উঃ} & \textbf{নং} & \textbf{উঃ} & \textbf{নং} & \textbf{উঃ} & \textbf{নং} & \textbf{উঃ} & \textbf{নং} & \textbf{উঃ} & \textbf{নং} & \textbf{উঃ} & \textbf{নং} & \textbf{উঃ} \\
\hline
\endhead
1 & ক & 2 & খ & 3 & গ & 4 & গ & 5 & খ & 6 & খ & 7 & ঘ \\
\hline
8 & ক & 9 & ক & 10 & ক & 11 & ঘ & 12 & ঘ & 13 & ক & 14 & গ \\
\hline
15 & গ & 16 & ক & 17 & ঘ & 18 & গ & 19 & ঘ & 20 & ক & 21 & গ \\
\hline
22 & খ & 23 & খ & 24 & ঘ & 25 & খ & 26 & খ & 27 & খ & 28 & ক \\
\hline
29 & খ & 30 & গ & 31 & ক & 32 & খ & 33 & গ & 34 & গ & 35 & ক \\
\hline
36 & ক & 37 & গ & 38 & ক & 39 & ঘ & 40 & গ & 41 & গ & 42 & গ \\
\hline
43 & ক & 44 & ঘ & 45 & ঘ & 46 & খ & 47 & খ & 48 & ঘ & 49 & ঘ \\
\hline
50 & খ & 51 & গ & 52 & গ & 53 & খ & 54 & ঘ & 55 & ঘ & 56 & গ \\
\hline
57 & গ & 58 & খ & 59 & খ & 60 & খ & 61 & গ & 62 & ক & 63 & ক \\
\hline
64 & গ & 65 & গ & 66 & খ & 67 & খ & 68 & ক & 69 & খ & 70 & গ \\
\hline
71 & খ & 72 & খ & -- & -- & -- & -- & -- & -- & -- & -- & -- & -- \\
\hline
\end{longtable}
\endgroup
\partheading{অংশ–৩ : Difficult, Critical ও Trap MCQ-এর ব্যাখ্যা}
\begin{infobox}
সব প্রশ্নের ব্যাখ্যা পুনরাবৃত্তি করা হয়নি। শুধু যেসব জায়গায় Trap আছে, Dry Run/Logic দরকার, বা একই Concept থেকে অন্য angle-এ প্রশ্ন আসার সম্ভাবনা বেশি—সেগুলো রাখা হয়েছে।
\end{infobox}
\begin{explainbox}{Topic-wise ২৮ — উত্তর: ঘ}
Interpreter সাধারণত Statement ধরে ধরে অনুবাদ/Execute করে। Error পাওয়া গেলে সেই জায়গাতেই কাজ থেমে যেতে পারে, তাই Error দ্রুত শনাক্ত করা যায়।

\textbf{এ Concept থেকে আর কীভাবে প্রশ্ন আসতে পারে?} Compiler-এর সাথে পার্থক্য, Error কখন দেখা যায়, Execution speed, Object Program তৈরি হয় কি না—এই angle-গুলো থেকেও প্রশ্ন হতে পারে।
\end{explainbox}

\begin{explainbox}{Topic-wise ৫৬ — উত্তর: গ}
Dry Run-এ Program বাস্তবে Run না করে হাতে ধাপে ধাপে Logic অনুসরণ করা হয়। Logical Error বোঝার জন্য এটি বিশেষ কার্যকর।

\textbf{এ Concept থেকে আর কীভাবে প্রশ্ন আসতে পারে?} Trace Table, Variable-এর পরবর্তী মান, Final Output, কোন Iteration-এ Condition False হবে—এভাবে প্রশ্ন আসতে পারে।
\end{explainbox}

\begin{explainbox}{Topic-wise ৬৩ — উত্তর: গ}
Flowchart-এ Input/Output-এর Standard Symbol হলো Parallelogram। Rectangle হলো Processing এবং Diamond হলো Decision।

\textbf{এ Concept থেকে আর কীভাবে প্রশ্ন আসতে পারে?} Symbol দেখে কাজ চেনা, কাজ দেখে Symbol চেনা, ভুল Symbol শনাক্ত—এই তিনভাবে প্রশ্ন ঘুরে আসতে পারে।
\end{explainbox}

\begin{explainbox}{Topic-wise ৮৩ — উত্তর: গ}
Python Identifier Letter বা underscore দিয়ে শুরু হতে পারে; Digit দিয়ে শুরু, hyphen বা Keyword ব্যবহার করা যায় না।

\textbf{এ Concept থেকে আর কীভাবে প্রশ্ন আসতে পারে?} Valid/invalid Variable Name, Case sensitivity, underscore, Keyword—এসব মিলিয়ে Statement-based MCQ হতে পারে।
\end{explainbox}

\begin{explainbox}{Topic-wise ১০৪ — উত্তর: ক}
\py{input()} সাধারণত String Return করে। Arithmetic দরকার হলে \py{int()} বা \py{float()} দিয়ে Convert করতে হয়।

\textbf{এ Concept থেকে আর কীভাবে প্রশ্ন আসতে পারে?} Input-এর Type, Conversion না করলে + কী করবে, int/float conversion—এসব নিয়ে Code Output প্রশ্ন হতে পারে।
\end{explainbox}

\begin{explainbox}{Topic-wise ১১৫ — উত্তর: ঘ}
Python-এ \textenglish{\ttfamily **} হলো Exponentiation Operator। তাই 2-এর power 3 = 8।

\textbf{এ Concept থেকে আর কীভাবে প্রশ্ন আসতে পারে?} Operator precedence, power-এর সাথে addition/multiplication, negative value—এই angle-এ প্রশ্ন আসতে পারে।
\end{explainbox}

\begin{explainbox}{Topic-wise ১২১ — উত্তর: খ}
Multiplication-এর precedence Addition-এর চেয়ে বেশি। তাই আগে 3×4 = 12, পরে 2+12 = 14।

\textbf{এ Concept থেকে আর কীভাবে প্রশ্ন আসতে পারে?}\par Parentheses যোগ করলে ফল, / // \% ** মিশিয়ে precedence এবং left-to-right evaluation—এগুলো practice করা দরকার।
\end{explainbox}

\begin{explainbox}{Topic-wise ১৩৪ — উত্তর: ঘ}
Python-এ Numeric 0 Boolean Context-এ False হিসেবে বিবেচিত হয়। তাই \py{if 0:} Block Execute হয় না।

\textbf{এ Concept থেকে আর কীভাবে প্রশ্ন আসতে পারে?} Non-zero Number, empty String/List, bool conversion—এসব থেকেও True/False নির্ণয়ের প্রশ্ন হতে পারে।
\end{explainbox}

\begin{explainbox}{Topic-wise ১৪৪ — উত্তর: ক}
\py{range(1, 6, 2)}-এ Start=1, Stop=6 (excluded), Step=2; তাই 1, 3, 5 পাওয়া যায়।

\textbf{এ Concept থেকে আর কীভাবে প্রশ্ন আসতে পারে?} Loop count, শেষ Value, Step পরিবর্তন, descending range—এগুলো থেকে ভিন্ন angle-এর প্রশ্ন খুব সম্ভব।
\end{explainbox}

\begin{explainbox}{Topic-wise ১৪৬ — উত্তর: ক}
\py{range(5,0,-1)} descending sequence তৈরি করে: 5,4,3,2,1। Stop 0 অন্তর্ভুক্ত হয় না।

\textbf{এ Concept থেকে আর কীভাবে প্রশ্ন আসতে পারে?} Negative step, invalid direction হলে empty range, কতবার Loop চলবে—এভাবে প্রশ্ন হতে পারে।
\end{explainbox}

\begin{explainbox}{Topic-wise ১৫২ — উত্তর: গ}
Loop-এর \py{else} সাধারণত Loop স্বাভাবিকভাবে শেষ হলে চলে; \py{break} দিয়ে বের হলে else চালানো হয় না।

\textbf{এ Concept থেকে আর কীভাবে প্রশ্ন আসতে পারে?} Search problem, break সহ/ছাড়া output, for-else ও while-else—এগুলো tricky MCQ হতে পারে।
\end{explainbox}

\begin{explainbox}{Topic-wise ১৭৫ — উত্তর: ঘ}
খালি Set-এর জন্য \py{set()} ব্যবহার করতে হয়; \textenglish{\ttfamily \{\}} আসলে Empty Dictionary।

\textbf{এ Concept থেকে আর কীভাবে প্রশ্ন আসতে পারে?} Set literal, Dictionary literal, type(), membership—এগুলো Trap হিসেবে আসতে পারে।
\end{explainbox}

\begin{explainbox}{Topic-wise ১৭৭ — উত্তর: গ}
Dictionary-তে \py{in} Operator সাধারণত Key-এর উপস্থিতি পরীক্ষা করে, Value-এর নয়।

\textbf{এ Concept থেকে আর কীভাবে প্রশ্ন আসতে পারে?} Key lookup, value membership, \py{keys()}, \py{values()}, \py{items()}—এগুলো থেকে প্রশ্ন হতে পারে।
\end{explainbox}

\begin{explainbox}{Topic-wise ১৯০ — উত্তর: ঘ}
কোনো Function explicit \py{return} না দিলে Python সাধারণত \py{None} Return করে। Screen-এ কিছু Print হওয়া Return Value নয়।

\textbf{এ Concept থেকে আর কীভাবে প্রশ্ন আসতে পারে?} print বনাম return, return-এর পরে Statement Execute হয় কি না, Function Call-এর Result Variable-এ রাখলে কী হয়—এভাবে প্রশ্ন হতে পারে।
\end{explainbox}

\begin{explainbox}{Topic-wise ১৯৩ — উত্তর: গ}
Function-এর ভিতরে Module-level Global Variable-এ Assignment/Rebinding করতে \py{global} Keyword প্রয়োজন হতে পারে। শুধু Value পড়তে সাধারণত লাগে না।

\textbf{এ Concept থেকে আর কীভাবে প্রশ্ন আসতে পারে?} Global read বনাম write, same-name local variable, output prediction—এসব গুরুত্বপূর্ণ।
\end{explainbox}

\begin{explainbox}{Board-style ১৬ — উত্তর: ক}
17 // 5 = 3 এবং 17 \% 5 = 2। তাই a+b = 5। একই প্রশ্নে Floor Division ও Modulus একসঙ্গে পরীক্ষা করা হয়েছে।

\textbf{এ Concept থেকে আর কীভাবে প্রশ্ন আসতে পারে?} Quotient/remainder, / বনাম //, negative operand থাকলে ফল—এসব angle-এ প্রশ্ন আসতে পারে।
\end{explainbox}

\begin{explainbox}{Board-style ২৫ — উত্তর: খ}
Negative step -2 হওয়ায় 5 থেকে 2 করে কমে 5,3,1 পাওয়া যায়; Stop 0 অন্তর্ভুক্ত নয়।

\textbf{এ Concept থেকে আর কীভাবে প্রশ্ন আসতে পারে?} Step -1/-2, Loop count, শেষ Value, Start-Stop direction ভুল হলে কী হয়—এসব থেকে প্রশ্ন হতে পারে।
\end{explainbox}

\begin{explainbox}{Board-style ২৮ — উত্তর: ক}
i == 3 হলে \py{break} Loop সম্পূর্ণ বন্ধ করে দেয়। তাই 1 ও 2 Print হওয়ার পর Loop শেষ।

\textbf{এ Concept থেকে আর কীভাবে প্রশ্ন আসতে পারে?} break-এর অবস্থান print-এর আগে/পরে বদলানো, break condition বদলানো, loop else—এইভাবে ঘুরিয়ে আসতে পারে।
\end{explainbox}

\begin{explainbox}{Board-style ২৯ — উত্তর: খ}
i == 3 হলে শুধু ওই Iteration Skip হয়; Loop বন্ধ হয় না। তাই 1,2,4,5 Print হয়।

\textbf{এ Concept থেকে আর কীভাবে প্রশ্ন আসতে পারে?} একাধিক continue, condition বদলানো, total/count-এর উপর continue-এর প্রভাব—এসব থেকে প্রশ্ন আসতে পারে।
\end{explainbox}

\begin{explainbox}{Board-style ৪০ — উত্তর: গ}
\py{set()} Empty Set তৈরি করে। Empty braces \textenglish{\ttfamily \{\}} Dictionary তৈরি করে—এটি খুব common trap।

\textbf{এ Concept থেকে আর কীভাবে প্রশ্ন আসতে পারে?} type(set()), type({}), duplicate removal, set indexing—এসব angle-এ প্রশ্ন হতে পারে।
\end{explainbox}

\begin{explainbox}{Board-style ৪৬ — উত্তর: খ}
Function-এর ভিতরের x = 20 একটি Local Variable; বাইরের x = 10 Global এবং অপরিবর্তিত থাকে। তাই Output 20 তারপর 10।

\textbf{এ Concept থেকে আর কীভাবে প্রশ্ন আসতে পারে?} global keyword যোগ করলে কী হবে, শুধু global x পড়লে কী হবে, parameter-এর নাম x হলে কী হবে—এসব থেকে প্রশ্ন আসতে পারে।
\end{explainbox}

\begin{explainbox}{Board-style ৫৫ — উত্তর: ঘ}
x-এর মান 1\(\rightarrow\)2\(\rightarrow\)4\(\rightarrow\)8\(\rightarrow\)16 হয়। 16 হলে x<10 False, তাই Loop শেষে x=16।

\textbf{এ Concept থেকে আর কীভাবে প্রশ্ন আসতে পারে?} কয়বার Loop চলল, print কোথায় বসালে কী Output, *=2-এর বদলে +=2—এভাবে প্রশ্ন হতে পারে।
\end{explainbox}

\begin{explainbox}{Board-style ৫৮ — উত্তর: খ}
nums[1:3]-এ Start index 1 অন্তর্ভুক্ত এবং Stop index 3 বাদ; তাই index 1 ও 2-এর Value [2,3] পাওয়া যায়।

\textbf{এ Concept থেকে আর কীভাবে প্রশ্ন আসতে পারে?} Negative slicing, omitted start/stop, step slicing, শেষ index excluded—এসব MCQ-তে আসতে পারে।
\end{explainbox}

\begin{explainbox}{Board-style ৬০ — উত্তর: খ}
Function-এর ভিতরে Global Variable শুধু পড়তে হলে সাধারণত global Keyword লাগে না; Assignment/Rebinding করতে হলে লাগে।

\textbf{এ Concept থেকে আর কীভাবে প্রশ্ন আসতে পারে?} Read বনাম modify, local shadowing, global declaration-এর অবস্থান—এগুলো critical question হতে পারে।
\end{explainbox}

\begin{explainbox}{Topic-wise ২০৪ — উত্তর: খ}
\py{x} ও \py{y} দুটিই String, তাই \py{+} এখানে যোগ না করে String Concatenation করে। ফলে \py{"5" + "2"} থেকে \py{"52"} পাওয়া যায়।

\textbf{Trap:} দেখতে Number হলেও Quote-এর ভিতরে থাকলে তা String।

\textbf{এ Concept থেকে আর কীভাবে প্রশ্ন আসতে পারে?} \py{int()} দিয়ে Conversion করলে Output কী হবে, একটি String ও একটি Integer সরাসরি \py{+} করলে কী Error হবে—এগুলো থেকেও প্রশ্ন আসতে পারে।
\end{explainbox}

\begin{explainbox}{Topic-wise ২০৮ — উত্তর: ক}
\py{range(5, 0, -2)}-এ Start = 5, Stop = 0 এবং Step = -2। তাই Value হবে \py{5, 3, 1}; Stop Value 0 অন্তর্ভুক্ত হয় না।

\textbf{এ Concept থেকে আর কীভাবে প্রশ্ন আসতে পারে?} Negative Step-এ Loop count, শেষ Value, Stop পরিবর্তন করলে Output—এই angle-গুলো গুরুত্বপূর্ণ।
\end{explainbox}

\begin{explainbox}{Topic-wise ২১২ — উত্তর: ঘ}
Outer Loop 2 বার এবং Inner Loop প্রতি Outer Iteration-এ 3 বার চলে। তাই \py{print("*")} মোট \(2 \times 3 = 6\) বার Execute হয়।

\textbf{Trap:} Nested Loop-এ Iteration count সাধারণত গুণ করতে হয়।

\textbf{এ Concept থেকে আর কীভাবে প্রশ্ন আসতে পারে?} Row–Column Pattern, মোট Output count, Inner range পরিবর্তন করলে মোট Iteration—এসব আসতে পারে।
\end{explainbox}

\begin{explainbox}{Topic-wise ২১৫ — উত্তর: খ}
Set Duplicate Element আলাদা করে রাখে না। তাই \py{\{1,2,2,3\}} কার্যত \py{\{1,2,3\}} হয় এবং Length = 3।

\textbf{এ Concept থেকে আর কীভাবে প্রশ্ন আসতে পারে?} List থেকে Set বানালে কতটি Unique Value থাকবে, Membership Test, Empty Set-এর Syntax—এসব থেকেও প্রশ্ন হতে পারে।
\end{explainbox}

\begin{explainbox}{Topic-wise ২১৯ — উত্তর: গ}
Function-এর ভিতরের \py{x = 20} একটি Local Variable; বাইরে থাকা \py{x = 10} Global Variable। Local Assignment বাইরের Global \py{x} পরিবর্তন করে না। তাই Output হবে 20, তারপর 10।

\textbf{এ Concept থেকে আর কীভাবে প্রশ্ন আসতে পারে?} \py{global x} ব্যবহার করলে Output কী হবে, Function-এর ভিতরে শুধু Global Variable Read করলে \py{global} দরকার কি না—এসব গুরুত্বপূর্ণ।
\end{explainbox}

\begin{explainbox}{Board-style ৬১ — উত্তর: গ}
\py{x} হলো String এবং \py{y} হলো Integer। Python-এ String ও Integer-কে সরাসরি \py{+} করা যায় না, তাই TypeError হবে।

\textbf{Trap:} \py{"10"} এবং \py{10} এক Data Type নয়।

\textbf{এ Concept থেকে আর কীভাবে প্রশ্ন আসতে পারে?} \py{int(x) + y}, \py{x + str(y)}, অথবা \py{input()}-এর Result নিয়ে একই ধরনের Output/Error প্রশ্ন আসতে পারে।
\end{explainbox}

\begin{explainbox}{Board-style ৬৪ — উত্তর: গ}
\py{i} শুরুতে 1 এবং Loop-এর ভিতরে \py{i}-এর Value পরিবর্তন করা হয়নি। তাই \py{i <= 5} সবসময় True থেকে যেতে পারে এবং Infinite Loop হতে পারে।

\textbf{এ Concept থেকে আর কীভাবে প্রশ্ন আসতে পারে?} সঠিক Update কোথায় বসবে, \py{i += 1} যোগ করলে কয়বার Loop চলবে, Condition পরিবর্তন করলে কী হবে—এসব angle থেকে প্রশ্ন আসতে পারে।
\end{explainbox}

\begin{explainbox}{Board-style ৬৫ — উত্তর: গ}
Outer Loop 3 বার এবং Inner Loop 2 বার চলে। তাই \py{count += 1} মোট \(3 \times 2 = 6\) বার Execute হয়।

\textbf{এ Concept থেকে আর কীভাবে প্রশ্ন আসতে পারে?} Nested Loop-এর Final count, Pattern-এর Row/Column, একটি Loop-এর range পরিবর্তন করলে Total Iteration—এসব আসতে পারে।
\end{explainbox}

\begin{explainbox}{Board-style ৬৮ — উত্তর: ক}
Dictionary-তে \py{in} Operator সাধারণত Key-এর Membership পরীক্ষা করে। এখানে \py{"a"} একটি Key, তাই Result \py{True}।

\textbf{Trap:} \py{10 in d} লিখলে Value দেখে True হবে—এমন ধরে নেওয়া ঠিক নয়।

\textbf{এ Concept থেকে আর কীভাবে প্রশ্ন আসতে পারে?} Key access, \py{keys()}, \py{values()}, Dictionary update বা Membership নিয়ে প্রশ্ন আসতে পারে।
\end{explainbox}

\begin{explainbox}{Board-style ৭২ — উত্তর: খ}
\py{show()} Function প্রথমে \py{Hi} Print করে। কিন্তু Function-এ কোনো \py{return} নেই, তাই এটি স্বয়ংক্রিয়ভাবে \py{None} Return করে। ফলে পরে \py{print(x)} থেকে \py{None} দেখা যায়।

\textbf{Trap:} \py{print()} এবং \py{return} একই কাজ করে না।

\textbf{এ Concept থেকে আর কীভাবে প্রশ্ন আসতে পারে?} Function-এ \py{return "Hi"} দিলে কী হবে, \py{print()} বাদ দিলে Result কী হবে, Return Value Variable-এ রাখলে কী পাওয়া যাবে—এসব থেকেও প্রশ্ন হতে পারে।
\end{explainbox}

\begin{rememberbox}
\textbf{Revision Rule:} গুরুত্বপূর্ণ Concept শুধু Definition হিসেবে নয়—Feature, Difference, Statement, Code Output, Dry Run, Trap ও Modification—সব angle থেকে Practice করবে। প্রশ্ন মুখস্থ নয়, Logic ধরাই মূল লক্ষ্য।
\end{rememberbox}

\RaggedRight
\small

